%!TeX root = ../incremental_SS_Translation.tex

\chapter{Śārīrasthāna 5:  On the enumeration of items in the body}


\section{Introduction}

%\texttt{%Draft text by Jan Gerris.
%    The Sanskrit text of this adhyāya is a lightly edited transcript of NAK 5-333,
%    waiting for full critical attention.}


\subsection{Literature}

Meulenbeld offered an annotated overview of this chapter and a
bibliography of earlier scholarship to 2002 and, in his notes, citations
of the parallel passages in the \CS.\fvolcite{IA}[249--254]{meul-hist}   

A translation of the  vulgate text of \SS\ 3.05 from the three
editions edited by Ācārya was published by
\citet[697--705]{zysk-1986}.\footnote{The three editions are
    \cite{susr-trikamji1,susr-trikamji2,susr-trikamji1945}.}

\newpage

\section{Translation}

\begin{translation}
  
  \begin{tt}
      \raggedright
      \bigskip
      \marginpar{Draft tr.\ from here}
      
\item[1]
Now, we shall explain the anatomical description of the bodyatha, 

\item[2]    vulgate only

\item[3]

The union of semen and female reproductive element located in the uterus, combined with the Self, the primary creative forces, and their modifications, is called the foetus. While it remains unconscious, the wind element divides it, the fire element digests (transforms) it, the water element moistens it, the earth element consolidates it, and the space element causes it to expand. When it is thus developed and endowed with hands, feet, tongue, nose, face, ears, and so forth, it then receives the designation "body." And that (body) has six parts: the four limbs, the trunk as the fifth, and the head as the sixth.

\item[4]
Beyond this point, the sub-limbs will be described: the head, belly, back, navel, forehead, nose, chin, and neck are each single.  The ears, eyes, eyebrows, temples, shoulders, cheeks, armpits, breasts, groins, testicles, sides, buttocks, knees, elbows, arms, and so forth, occur in pairs; the twenty fingers and the channels will be described. Thus, this division of sub-limbs has been stated, along with the membranes, the elements, and the carriers of waste  covered by the skin

\item[5, 6 and 7]  vulgate only

\item[8]

There are seven receptacles for wind, bile, phlegm, undigested food, digested food, metabolic fire, and urine.  In women, the uterus is the eighth. 

\item[9]

The intestines of men measure three and a half vyāmas (fathoms), while those of women are less by half a vyāma. 

\item[10]

Women have these same ones and three others. 
Two are in the breasts and another is below, carrying the menstrual blood. 

\item[11]
There are sixteen large tendons; of these, four are in the feet, and an equal number are in the hands, neck, and back. Among those, the tendons located in the hands and feet have nine branches each; [likewise] those connected to the neck and heart, those connected to the penis, blood, and back that go downwards, and those of the buttocks, thighs, and calves. 

\item[12]
There are four types of plexuses each—muscular, vascular, ligamentous, and bony—and these are situated in the wrists and ankles. By these, this body is overspread. 

\item[13]
There are six brush-like structures; they are in the hands, feet, neck, and penis. 

\item[14]
There are four great muscular cords for the purpose of binding the muscles; two are external and two are internal on either side of the spinal column. 

\item[15]
There are seven sutures; they are distributed as five in the head and one each in the tongue and the penis; these should be avoided by the surgical instrument. 

\item[16]
There are fourteen bone-unions; of these, three are located in the ankle, knee, and groin.
By this, the other leg and the two arms are explained. 
There is one in each of the sacrum and the head. 

\item[17]
There are exactly fourteen sutures. 
Those unions from which they are released are indeed eighteen according to some. 

\item[18]
The followers of the Vedas state that there are three hundred and sixty bones. In the surgical treatises, there are only three hundred. Of these, one hundred and twenty bones are in the extremities.  One hundred and seventeen are in the pelvis, sides, back, and chest; and sixty-three in the neck and the regions above. 

\item[19]
In each toe of the foot there are three; these make fifteen, located in the sole, the cluster, and the ankle. Ten are in the heel, two in the lower leg, one in the knee, and one in the thigh Thus, there are thirty bones in one leg; by this, the other leg and the two arms are explained.  There are five in the pelvis.  Of those, four are in the anus, genitals, and hips, and one is supported by the sacrum.  There are thirty-six in one side, and just as many in the second. There are thirty-two in the back and eight in the chest. Two are known as collar-bones; there are nine in the neck, four in the windpipe.  Two are known as the jaws, thirty-two are the teeth, three in the nose, one in the palate, and one each in the cheeks, ears, and temples, and six in the head. 

\item[20]
These bones are of five types. They are as follows: They are known as the flat, the rucaka, the cartilaginous, the curved, and the reed-like bones. Among them, the flat bones are located in the knees, hips, cheeks, palate, and the cranium.  The teeth, however, are the rucaka bones. The cartilaginous bones are found in the nose, throat, ears, and the eye sockets. The curved bones are in the feet, hands, neck, and back. 

\item[21]
And there is a verse on this subject: just as the core essence of a tree remains situated within, so too the body of living beings is supported by the bones. 

\item[22]
Therefore, even when the skin and flesh of embodied beings have long since perished, the bones do not decay, as they are considered the core essence. 

\item[23]
The muscles are bound here by the vessels and tendons; taking the bones as their support, they do not waste away or fall off. 

\item[24]
Joints are of two kinds: those that are movable and those that are immovable. 
\item[25]   vulgate only

\item[26]
They are counted as two hundred and ten; of these, sixty-eight are in the extremities, fifty-nine in the trunk, and eighty-three in the neck and the region above it. In each toe there are three joints, and two in the big toe; these fourteen, along with one each in the ankle, knee, and groin, make seventeen in one leg.  By this, the other leg and the two arms are also explained. There are three in the pelvic bones and twenty-four in the spinal column.  There are the same number in the two sides. There are eight in the chest.  There are the same number in the neck; eighteen are in the windpipe and the vessels connected to the heart and kloma. There are as many as the teeth in the roots of the teeth, one in each. There is one each in the thyroid and the nose Two are in the circular eyelids, resting on the eyes. 
There is one each in the cheeks, ears, and temples.  There are two jaw joints, two above the eyebrows and temples, five in the skull bones, and one in the crown of the head. 

\item[27]
These joints are of eight types. They are as follows: Kora (hinge), Ulūkhala (mortar), Sāmudga (box), Pratara (flat), Tunna-sevanī (suture), Vāyasa-tuṇḍa (crow's beak), Maṇḍala (circular), and Śaṅkhāvarta. Among them, the Kora joints are in the fingers, wrists, ankles, knees, and elbows The Ulūkhala joints are in the armpits, groins, and teeth. The Sāmudga joints are in the shoulders, hips, anus, genitals, and buttocks The Pratara joints are in the neck and spinal column. The Tunna-sevanī joints are in the skull and pelvic bones.  The Vāyasa-tuṇḍa joints are on both sides of the jaws.  The Maṇḍala joints are in the throat, heart, eyes, eyelids, and the kloma-vessels.  The Śaṅkhāvarta joints are in the ears and the cranial vascular junctions.  Their shapes are generally explained by their very names. 

\item[28]
And there is a verse on this subject: only these bone-joints have been fully described; the number of joints of the muscles, ligaments, and veins is not known.  


\item[29]
There are nine hundred ligaments. Of these, six hundred are in the extremities.  Two hundred and thirty are in the trunk, and seventy are in the neck and above.  In each individual toe of the foot, there are six, totalling thirty. The same number is found in the soles, the ankles, and the ankle-joints, and the same number in the lower leg.  Ten are in the knee. Forty are in the thighs. Ten are in the groin. Thus, there are one hundred and fifty in a single leg.  By this, the other leg and the two arms are explained. Sixty are in the hips, while eighty are in the back. There are sixty in the two sides.  Thirty are in the chest. Thirty-six are in the neck. Thirty-four are in the head.  Thus, nine hundred ligaments have been explained. Among the large ligaments, the technical term is "tendon" (kaṇḍarā). 


\item[30]
And there are verses on this subject: Understand from me that ligaments are of four types; they are not otherwise. Branching, circular, porous, and broad

\item[31]
Furthermore, the branching ones are in the limbs and in all the joints All the circular ones should be known by experts here as tendons. 
\item[32]
Broad and porous [ligaments] are found at the ends of the stomach and intestines, as well as in the sides, the chest, the back, and the head.

\item[33]
Just as a boat made of planks, when joined with many fastenings, becomes capable of bearing heavy loads because it relies on firm joints, 

\item[34]
in this very same way, all the joints that are recorded in this body are bound by many ligaments and are therefore known to be capable of bearing weight. 


\item[35]
For injuries to the bones, muscles, veins, or joints do not harm embodied beings as much as an injury to the ligaments does. 

\item[36]
The man who understands the ligaments, both external and internal, is able to extract a deeply embedded foreign object from the body. 

\item[37]
There are five hundred muscles; of those, four hundred are in the extremities, sixty are in the trunk, and forty are in the neck and above. In each toe of the foot, there are three [muscles] each; these make fifteen. There are just as many situated in the Kūrca.Ten are in the ankle and the sole.  There are twenty between the ankle and the knee.  Five are in the knee. Ten are in the thighs.  Ten are in the groin.  There are one hundred in one leg; by this, the other leg and the two arms are also explained. One is in the penis and the perineum; Two others are in the scrotum. Three are in the anus.  And five each in the two buttocks. Two are at the head of the bladder. Seven are in the abdomen and one in the navel.  Five long ones each are situated in the upper part of the back. Ten are in the two sides. Two are in the chest. Seven are all around the collarbones Two are in the heart and the stomach. One is inside the trunk. And four are in the neck. One each is in the throat and the larynx. Two are in the palate.  Six are in the tongue. Two are  n the lips. Two are in the nose. Four each are in the two cheeks. Two are in the eyes. Two are in the ears. Four are in the forehead. Six are in the head.These are the five hundred muscles of men. 

\item[38]  vulgate only

\item[39]
In women, there are twenty more.  Ten of those are in the breasts, five in each. 
Their full growth occurs during youth.  Two are located at the vaginal orifice
Two are circular and support the foetus.  Three serve as the entry points for semen and menstrual blood. There are exactly three in the uterus where the foetus resides. 

\item[40 \& 41]   vulgate only

\item[42]
The classification of vital points, veins, arteries, and channels is described elsewhere. 

\item[43]
And there is a verse on this subject: The female reproductive organ is famously described as being shaped like the navel of a conch shell and possessing three spiral turns. 
The bed of the foetus is situated within its third turn. 

\item[44]
The shape of the uterus resembles the mouth of a Rohita fish, and the wise know the bed of the foetus to be of that same form. 

\item[45]
The foetus lies in a crouched position within the mother's womb, facing her; by nature, it moves toward the birth canal with its head first for delivery. T

\item[46]
This detailed determination of the body's parts, extending to the skin, is not described in any other branches of study except for surgical knowledge.

\item[47]
Therefore, a surgeon who desires certain knowledge must thoroughly examine the determination of the body's parts by preparing and dissecting a cadaver

\item[48]
For when that which is perceived through direct observation and that which is perceived through the scriptures are combined, both together greatly increase one's knowledge. 

\item[49]
Therefore, one should take the body of a person that has all its limbs intact and was not afflicted by a long-standing disease, place it in a cage, and, having bound it with Muñja grass and Balvaja grass, let it decompose in a hidden spot in a river. When it is fully decomposed, one should remove the body and, having laid it out, should slowly scrub it with any one of a brush of Uśīra grass, bamboo, or bark, and personally observe with the eyes all the internal and external parts and specific organs as they have been described. 
\item[50]
There are two verses on this.  The wise know that the extremely subtle cannot be seen with the physical eye; it is perceived only with the eyes of knowledge or the eyes of asceticism


\item[51]
An expert is one who has observed the facts both in the actual body and in the scriptures; having dispelled doubt through both observation and study, he begins his medical practice. 

Thus ends the fifth chapter in the Section on the Body
\end{tt}
\end{translation}